\documentclass[11pt]{article}

\topmargin=-0.5in
\oddsidemargin=0in
\evensidemargin=0in
\textwidth=6.5in
\textheight=9.0in

\usepackage{setspace}
\usepackage{graphicx}

\usepackage{amsmath}
\usepackage{amssymb}
\usepackage{amsthm}
\usepackage{url}

\usepackage{setspace}

\input{Qcircuit}

% \noindent {\sc Slide}: Jones-Gomard-Sestoft book cover.

% \mbox{} \\ \hspace*{1cm}
%    [Our experiments confirm our theory] \\

\title{Quantum Algorithms \\
       Clifford+T is Universal}
\author{Jens Palsberg}
\date{Mar 5, 2025}

% Title: Towards a universal quantum programming language.
% Goal: 2,500 words.
% Almost every quantum computer right now has its own programming language. 
% How may this change for the future quantum computers and what are 
% the unique challenges in making a universal quantum computing language?
%
% Role model: http://xrds.acm.org/current-issue.cfm

\def\({\begin{eqnarray*}}
\def\){\end{eqnarray*}}

\newtheorem{theorem}{Theorem}
\newtheorem{corollary}[theorem]{Corollary}
\newtheorem{lemma}[theorem]{Lemma}

\def\phi{\varphi}
\def\imp{\rightarrow}
\def\implies{\Rightarrow}

\def\ket#1{|#1\rangle}
\def\bra#1{\langle#1|}
\def\inner#1#2{\langle{#1}|{#2}\rangle}
\def\outer#1#2{\ket{#1}\bra{#2}}
\def\eigenvalues#1{\mbox{\sf Eigenvalues}(#1)}

\def\iff{\Longleftrightarrow}

\begin{document}
\maketitle

% Example
% Story
% Data
% Quote
% Peer Testimony

\newpage

\rightline{Quantum Algorithms, by Jens Palsberg}
\rightline{Clifford+T is Universal; 100 minutes; Mar 5, 2025} 

\vspace*{-1.2cm}

\section*{Outline}

\paragraph{Hook:} 
A quantum computer implements a universal gate set.
Intuitively, this means that the computer can implement any unitary. 
But technically, it means that the computer can approximate any unitary.
What is behind this idea and how can we prove that 
a gate set can approximate any unitary?

\paragraph{Purpose:} Persuade you that Clifford+T is universal.

\paragraph{Preview:} \mbox{} \\
\hspace*{1cm}
1. We will state the theorem and devise a proof strategy.
\\
\hspace*{1cm}
2. We will first prove a version of the theorem that works with $R_z$ gates.
\\
\hspace*{1cm}
3. We will prove the full theorem by approximating the $R_z$ gates.

\paragraph{Transition to Body:} 
Let us begin with a look at what we are trying to prove.

\paragraph{Main Point 1:}
We will state the theorem and devise a proof strategy. \\
\hspace*{1cm}
   [We state theorem] \\ 
\hspace*{1cm}
   [We will first prove a version of the theorem that uses an easy gate set] \\
\hspace*{1cm}
   [We state some known lemmas]

\paragraph{Transition to MP2:}
Now we move on to a proof for an easy gate set.

\paragraph{Main Point 2:}
We will first prove a version of the theorem that works with $R_z$ gates. \\
\hspace*{1cm}
   [We prove the simple version of the theorem] \\
\hspace*{1cm}
   [We prove a key lemma] \\
\hspace*{1cm}
   [We need exponentially many gates to implement an $n$-qubit unitary in the worst case] 

\paragraph{Transition to MP3:}
Now we can finish the proof by focusing on approximation.

\paragraph{Main Point 3:}
We will prove the full theorem by approximating the $R_z$ gates. \\
\hspace*{1cm}
   [We define approximation of unitaries] \\
\hspace*{1cm}
   [We state some known lemmas about the Hilbert-Schmidt distance] \\
\hspace*{1cm}
   [We prove the full version of the theorem] 

\paragraph{Transition to Close:} 
So there you have it.

\paragraph{Review:}
The key property of a universal gate set is that 
it can approximate any unitary.
Our proof strategy is to first prove the theorem for an easy gate set and then
extend to $\{C(X), H, T\}$.

\paragraph{Strong finish:}
Quantum computing has probabilities built in and
the idea of a universal gate set has approximation built in.
Quantum algorithms have no requirement to be exact; 
rather, they can try to get close.

\paragraph{Call to action:}
Implement the procedure for compiling any unitary into the easy gate set.

\newpage

\section*{Detailed presentation}

\paragraph{Hook:} 
A quantum computer implements a universal gate set.
Intuitively, this means that the computer can implement any unitary. 
But technically, it means that the computer can approximate any unitary.
What is behind this idea and how can we prove that 
a gate set can approximate any unitary?

\paragraph{Purpose:} Persuade you that Clifford+T is universal.

\paragraph{Preview:} \mbox{} \\
\hspace*{1cm}
1. We will state the theorem and devise a proof strategy.
\\
\hspace*{1cm}
2. We will first prove a version of the theorem that works with $R_z$ gates.
\\
\hspace*{1cm}
3. We will prove the full theorem by approximating the $R_z$ gates.

\paragraph{Transition to Body:} 
Let us begin with a look at what we are trying to prove.

\paragraph{Main Point 1:}
We will state the theorem and devise a proof strategy. \\
\mbox{} \\ \hspace*{1cm}
   [We state theorem] \\ 
Let $d$ denote the Hilbert-Schmidt distance.

\bigskip

We will use the term circuit to mean a product of unitaries.
A key idea is that we can restrict those unitaries to come from
a particular gate set.  

We say that a gate set is universal if,
for any unitary $U$ and any $\epsilon > 0$, 
we can use the gate set to define a circuit $C$
such that $d(U,C) < \epsilon$.
The following theorem states that $\{C(X), H, T\}$ is universal.

\bigskip

\noindent
{\bf Theorem}.  For any unitary $U$ and any $\epsilon > 0$, 
there exists a circuit $C$ built from $\{C(X), H, T\}$,
such that $d(U,C) < \epsilon$.

\mbox{} \\ \hspace*{1cm}
   [We will first prove a version of the theorem that uses an easy gate set] \\
Let us state an easier version that has no need for approximation.

\begin{lemma} 
For any unitary $U$,
there exists a circuit $C$ built from 
2-qubit gates and
$\{H, S, S^\dagger, R_z(\theta)\}$, \\
such that $U = C$.
\label{lem:decomposition-to-an-easy-gate-set}
\end{lemma}

Our proof strategy is to first prove 
Lemma~\ref{lem:decomposition-to-an-easy-gate-set} 
and then approximate the $R_z(\theta)$ gates.

\mbox{} \\ \hspace*{1cm}
   [We state some known lemmas] \\
Let us define three rotation matrices:
\[
\begin{array}{rclclcl}
R_x(\theta) &=& 
e^{-i\theta X/2}
&=&
\cos(\theta/2)\ I - i\sin(\theta/2)\ X
&=&
\left(
\begin{array}{cc}
\cos(\theta/2)    & -i \sin(\theta/2) \\
-i \sin(\theta/2) & \cos(\theta/2)
\end{array}
\right)
\\
R_y(\theta) &=&
e^{-i\theta Y/2}
&=&
\cos(\theta/2)\ I - i\sin(\theta/2)\ Y
&=&
\left(
\begin{array}{cc}
\cos(\theta/2)    & -\sin(\theta/2) \\
\sin(\theta/2) & \cos(\theta/2)
\end{array}
\right)
\\
R_z(\theta) &=&
e^{-i\theta Z/2}
&=&
\cos(\theta/2)\ I - i\sin(\theta/2)\ Z
&=&
\left(
\begin{array}{cc}
e^{-i\theta/2} & 0 \\
0              & e^{i\theta/2}
\end{array}
\right)
\end{array}
\]

\newpage

Our next lemma states the Cosine-Sine decomposition,
also known as the CS-decomposition.

\begin{lemma} {\bf (Paige and Wei 1994)}
For one n-qubit gate $U$, there exist $2^{n-1}$ 1-qubit gates $R_y(\theta_0)$, $R_y(\theta_1)$, $\cdots$, and $R_y(\theta_{2^{n-1}-1})$, four (n-1)-qubit gates $P_0$, $P_1$, $Q_0$, and $Q_1$, and three n-qubit gates $P = \outer{0}{0} \otimes P_{0} + \outer{1}{1} \otimes P_{1}$, $R = \sum^{2^{n-1}-1}_{i=0}R_y(\theta_i) \otimes \outer{i}{i}$, and $Q= \outer{0}{0} \otimes Q_{0} + \outer{1}{1} \otimes Q_{1}$, 
such that
\[ 
\Qcircuit @C=0.5em @R=0.4em {
  & \qw & \multigate{2}{U} & \qw & &&  & & \qw & \gate{} \qwx[2]& \gate{R}      &\gate{} \qwx[2]&\qw\\
 &  &                &     & &=& &   &    &          &                 &               &   \\
 & \qw / & \ghost{U}        & \qw &  && & & \qw / & \gate{Q}       & \gate{} \qwx[-2]&\gate{P}       &\qw\\
}
\]
\label{cosinesine}
\end{lemma}

Our next lemma says that we can implement an $R_y(\alpha)$ via use of 
an $R_z(-\alpha)$ gate.

\begin{lemma} {\bf (Shende, Markov, and Bullock 2004)}
For one 1-qubit gate $R_y(\alpha)$, there exists one 1-qubit gate $R_z(-\alpha)$, such that
\[ 
\Qcircuit @C=0.5em @R=0.4em {
 & \gate{R_y(\alpha)} & \qw & & = & &  & \gate{S} & \gate{H} & \gate{R_z(-\alpha)}      &\gate{H}    & \gate{S^\dag} & \qw \\ 
}
\]
\label{ryrz}
\end{lemma}

\begin{proof}
We calculate:
%
\(
S^\dagger\ H\ R_z(-\alpha)\ H\ S
&=&
S^\dagger\ H\ [\cos(-\alpha/2)\ I - i\sin(-\alpha/2)\ Z]\ H\ S
\\
&=&
S^\dagger\ H\ [\cos(\alpha/2)\ I + i\sin(\alpha/2)\ Z]\ H\ S
\\
&=&
S^\dagger\ [\cos(\alpha/2)\ I + i\sin(\alpha/2)\ X]\ S
\\
&=&
\cos(\alpha/2)\ I + i\sin(\alpha/2)\ (-Y)
\\
&=&
\cos(\alpha/2)\ I - i\sin(\alpha/2)\ Y
\\
&=&
R_y(\alpha)
\)
In the third step, we use $H H = I$ and $H Z H = X$.
In the fourth step, we use $S^\dagger S = I$ and $S^\dagger X S = - Y$.
\end{proof}

The following lemma is interesting in itself, yet for our purposes,
the main point is that it works well together with CS-composition.

\begin{lemma} {\bf (Shende, Bullock, and Markov 2005)}
For two (n-1)-qubit gates $U_{0}$ and $U_{1}$ and one n-qubit gate $U = \outer{0}{0} \otimes U_{0} + \outer{1}{1} \otimes U_{1}$, there exist $2^{n-1}$ 1-qubit gates $R_z(\alpha_0)$, $R_z(\alpha_1)$, $\cdots$, and $R_z(\alpha_{2^{n-1}-1})$, two (n-1)-qubit gates $P$ and $Q$, and one n-qubit gate $R = \sum^{2^{n-1}-1}_{i=0}R_{z}(\alpha_i) \otimes \outer{i}{i}$, such that
\[ 
\Qcircuit @C=0.5em @R=0.4em {
& \qw & \gate{} \qwx[2] & \qw &  && &  & \qw & \qw       & \gate{R}      &\qw            &\qw\\  
& & &     & & = & &  &             &        &         &               &   \\
& \qw / & \gate{U}       & \qw &  && &  & \qw / & \gate{Q}       & \gate{} \qwx[-2]&\gate{P}       &\qw\\
}
\]
\label{demultiplexing}
\end{lemma}

\paragraph{Transition to MP2:}
Now we move on to a proof for an easy gate set.

\newpage

\paragraph{Main Point 2:}
We will first prove a version of the theorem that works with $R_z$ gates. \\
\mbox{} \\ \hspace*{1cm}
   [We prove the simple version of the theorem] \\
We state Lemma 1 again.

\bigskip

\noindent
{\bf Lemma 1.} 
{\em 
For any unitary $U$,
there exists a circuit $C$ built from 
2-qubit gates and
$\{H, S, S^\dagger, R_z(\theta)\}$, \\
such that $U = C$.}

\begin{proof}
For a $n$-qubit gate $U$, we calculate as follows.

\begin{eqnarray}
\Qcircuit @C=0.5em @R=1em {
& \qw & \multigate{2}{U}  & \qw \\
& \qw & \ghost{U} & \qw \\
& / \qw & \ghost{U}  & \qw 
}
& \;\;\;\;\;\;\; &
\Qcircuit @C=0.5em @R=1em {
& \qw & \gate{} \qwx[1] & \gate{R}  & \gate{} \qwx[1] & \qw \\
\lstick{= \;\;\;\; } & \qw & \multigate{1}{Q} &  \gate{} \qwx[-1] & \multigate{1}{P} & \qw \\
& / \qw & \ghost{Q}  &  \gate{} \qwx[-1] & \ghost{P}  & \qw 
}\nonumber
\\[0.4cm] &&
\Qcircuit @C=0.5em @R=1em {
& \qw & \gate{} \qwx[1] & \gate{S} & \gate{H} & \gate{R_1} & \gate{H} &\gate{S^\dag} & \gate{} \qwx[1] & \qw \\
\lstick{= \;\;\;\; } & \qw & \multigate{1}{Q} & \qw & \qw &  \gate{} \qwx[-1] & \qw & \qw & \multigate{1}{P} & \qw \\
& / \qw & \ghost{Q} & \qw & \qw &  \gate{} \qwx[-1] & \qw & \qw & \ghost{P}  & \qw 
}\nonumber
\\[0.4cm] &&
\Qcircuit @C=0.5em @R=1em {
& \qw & \gate{} \qwx[1] & \gate{S} & \gate{H} & \gate{R_1} & \gate{H} &\gate{S^\dag} & \qw & \gate{R_0} \qwx[1] & \qw & \qw \\
\lstick{= \;\;\;\; } & \qw & \multigate{1}{Q} & \qw & \qw &  \gate{} \qwx[-1] & \qw & \qw & \multigate{1}{W} & \gate{} \qwx[-1] & \multigate{1}{V} & \qw \\
& / \qw & \ghost{Q} & \qw & \qw &  \gate{} \qwx[-1] & \qw & \qw & \ghost{W} & \gate{} \qwx[-1] & \ghost{V}  & \qw 
}\nonumber
\\[0.4cm] &&
\Qcircuit @C=0.5em @R=1em {
& \qw & \qw & \gate{R_2} \qwx[1] & \qw & \gate{S} & \gate{H} & \gate{R_1} & \gate{H} &\gate{S^\dag} & \qw & \gate{R_0} \qwx[1] & \qw & \qw \\
\lstick{= \;\;\;\; } & \qw & \multigate{1}{N} & \gate{} \qwx[-1] & \multigate{1}{M} & \qw & \qw &  \gate{} \qwx[-1] & \qw & \qw & \multigate{1}{W} & \gate{} \qwx[-1] & \multigate{1}{V} & \qw \\
& / \qw & \ghost{N} & \gate{} \qwx[-1] & \ghost{M} & \qw & \qw &  \gate{} \qwx[-1] & \qw & \qw & \ghost{W} & \gate{} \qwx[-1] & \ghost{V}  & \qw 
}
\nonumber
% \label{eq:sec2}
\end{eqnarray}

In the first step, according to Lemma~\ref{cosinesine}, for every $n$-qubit gate $U$, there exist $2^{n-1}$ 1-qubit gates $R_y(\theta_0)$, $R_y(\theta_1)$, $\cdots$, and $R_y(\theta_{2^{n-1}-1})$, four (n-1)-qubit gates $P_0$, $P_1$, $Q_0$, and $Q_1$, and three n-qubit gates $P = \outer{0}{0} \otimes P_{0} + \outer{1}{1} \otimes P_{1}$, $R = \sum^{2^{n-1}-1}_{i=0}R_y(\theta_i) \otimes \outer{i}{i}$, and $Q= \outer{0}{0} \otimes Q_{0} + \outer{1}{1} \otimes Q_{1}$, 
such that the first step is valid. 

In the second step, according to Lemma~\ref{ryrz}, for $R = \sum^{2^{n-1}-1}_{i=0}R_y(\theta_i) \otimes \outer{i}{i}$ gate, there exists one $n$-qubit gate $R_1 = \sum^{2^{n-1}-1}_{i=0}R_z(-\theta_i) \otimes \outer{i}{i}$, such that 
\begin{eqnarray*}
R & = & \sum^{2^{n-1}-1}_{i=0}R_y(\theta_i) \otimes \outer{i}{i} \\
& = & \sum^{2^{n-1}-1}_{i=0} ( S^\dag \cdot H \cdot R_z(-\theta_i) \cdot H \cdot S )\otimes \outer{i}{i} \\
& = & \sum^{2^{n-1}-1}_{i=0} ( S^\dag \cdot H )\otimes \outer{i}{i} \cdot  \sum^{2^{n-1}-1}_{i=0}   R_z(-\theta_i) \otimes \outer{i}{i}  \cdot \sum^{2^{n-1}-1}_{i=0} ( H \cdot S )\otimes \outer{i}{i} \\
& = & (S^\dag \otimes I^{n-1}) \cdot (H \otimes I^{n-1}) \cdot R_1 \cdot (H \otimes I^{n-1}) \cdot (S \otimes I^{n-1})
\end{eqnarray*}

From the above, the second step is valid. 

In the third step, according to Lemma~\ref{demultiplexing}, for $P = \outer{0}{0} \otimes P_{0} + \outer{1}{1} \otimes P_{1} $ gate, there exist $2^{n-1}$ 1-qubit gates $R_z(\alpha_0)$, $R_z(\alpha_1)$, $\cdots$, and $R_z(\alpha_{2^{n-1}-1})$, two (n-1)-qubit gates $V$ and $W$, and one n-qubit gate $R_0 = \sum^{2^{n-1}-1}_{i=0}R_{z}(\alpha_i) \otimes \outer{i}{i}$, such that the third step is valid. 

In the fourth step, according to Lemma~\ref{demultiplexing}, for $Q = \outer{0}{0} \otimes Q_{0} + \outer{1}{1} \otimes Q_{1} $ gate, there exist $2^{n-1}$ 1-qubit gates $R_z(\beta_0)$, $R_z(\beta_1)$, $\cdots$, and $R_z(\beta_{2^{n-1}-1})$, two (n-1)-qubit gates $M$ and $N$, and one n-qubit gate $R_2 = \sum^{2^{n-1}-1}_{i=0}R_{z}(\beta_i) \otimes \outer{i}{i}$, such that the fourth step is valid.

The above calculation shows how to recursively decompose any unitary $U$.
\end{proof}
 
\mbox{} \\ \hspace*{1cm}
   [We prove a key lemma] \\
We first state a lemma that gets us from control on $(n-1)$ qubits to
control on $(n-2)$ qubits for the case of control of $R_z(\alpha)$ gates.

\begin{lemma} {\bf (M{\"{o}}tt{\"{o}}nen et al.\ 2004)}
For $2^{n-1}$ 1-qubit gates $R_z(\alpha_0)$, $R_z(\alpha_1)$, $\cdots$, and $R_z(\alpha_{2^{n-1}-1})$ and one n-qubit gate $R= \sum^{2^{n-1}-1}_{i=0}R_{z}(\alpha_i) \otimes \outer{i}{i}$, there exist $2^{n-2}$ 1-qubit gates $R_z(\beta_0)$, $R_z(\beta_1)$, $\cdots$, and $R_z(\beta_{2^{n-2}-1})$, $2^{n-2}$ 1-qubit gates $R_z(\gamma_0)$, $R_z(\gamma_1)$, $\cdots$, and $R_z(\gamma_{2^{n-2}-1})$, and two (n-1)-qubit gates $R_0 = \sum^{2^{n-2}-1}_{i=0}R_{z}(\beta_i) \otimes \outer{i}{i}$ and $R_1 = \sum^{2^{n-2}-1}_{i=0}R_{z}(\gamma_i) \otimes \outer{i}{i}$, such that
\[ 
\Qcircuit @C=0.5em @R=0.4em {
  & \qw & \gate{R} \qwx[1] & \qw & &&  & & \qw & \targ & \gate{R_0}      & \targ & \gate{R_1} & \qw\\
 &  \qw &   \gate{} \qwx[1]     &  \qw   & &=& &   &   \qw &     \ctrl{-1}     &  \qw               &        \ctrl{-1}       & \qw  & \qw \\
 & \qw / & \gate{}        & \qw &  && & & \qw / & \qw       & \gate{} \qwx[-2] & \qw  &\gate{} \qwx[-2]    &\qw\\
}
\]
\label{rzrz}
\end{lemma}

Now we can count the number of 2-qubit gates needed to implement 
multi-control $R_z(\alpha)$ gates.

\begin{lemma}
For $2^{n-1}$ 1-qubit gates $R_z(\alpha_0)$, $R_z(\alpha_1)$, $\cdots$, and $R_z(\alpha_{2^{n-1}-1})$ and one n-qubit gate $R= \sum^{2^{n-1}-1}_{i=0}R_{z}(\alpha_i) \otimes \outer{i}{i}$, we can implement the $R$ gate by using $(3 \times 2^{n-2} - 2)$ 2-qubit gates. 
\label{rzcount}
\end{lemma}

\begin{proof}
According to Lemma~\ref{rzrz}, for every n-qubit gate in the form of $R = \sum^{2^{n-1}-1}_{i=0}R_{z}(\alpha_i) \otimes \outer{i}{i}$, we can implement it by using two (n-1)-qubit gates $R_0 = \sum^{2^{n-2}-1}_{i=0}R_{z}(\beta_i) \otimes \outer{i}{i}$ and $R_1 = \sum^{2^{n-2}-1}_{i=0}R_{z}(\gamma_i) \otimes \outer{i}{i}$ and two $C(X)$ gates. From this, if we use $\mathcal{R}(n)$ to denote the 2-qubit gate count for an n-qubit gate in the form of $R = \sum^{2^{n-1}-1}_{i=0}R_{z}(\alpha_i) \otimes \outer{i}{i}$,  we know that 
\begin{eqnarray}
\mathcal{R}(n) & = &  2 \times \mathcal{R}(n-1) + 2 
\label{eq:-r-n} \\
\mathcal{R}(2) & = & 1
\label{eq:-r-2} 
\end{eqnarray}

Now we can do an iterative calculation that uses
Equations~(\ref{eq:-r-n})--(\ref{eq:-r-2}):
\begin{eqnarray}
\mathcal{R}(n) & = &  2 \times \mathcal{R}(n-1) + 2 \nonumber \\ 
& = & 2 \times (2 \times \mathcal{R}(n-2) + 2) + 2 \nonumber \\ 
& = & 2^2 \times \mathcal{R}(n-2) + 2^2 + 2 \nonumber \\ 
&& \cdots \nonumber \\  
& = & 2^{n-2} \times \mathcal{R}(2) + \sum^{n-2}_{i=1}2^i \nonumber \\ 
& = & 2^{n-2} + 2^{n-1} - 2 \nonumber \\ 
& = & 3 \times 2^{n-2} - 2 \nonumber
\end{eqnarray}

We conclude that we can implement the $R$ gate by $(3 \times 2^{n-2} - 2)$ 2-qubit gates. 
\end{proof}


\mbox{} \\ \hspace*{1cm}
   [We need exponentially many gates to implement an $n$-qubit unitary in the worst case] \\
From the above, we can implement every n-qubit gate $U$ by using 
four $(n-1)$-qubit gates $V$, $W$, $M$, and $N$ and 
three n-qubit controlled-rotation gates $R_0$, $R_1$, and $R_2$. 
If we use $\mathcal{U}(n)$ to denote the 2-qubit gate count for 
an general $n$-qubit gate and $\mathcal{R}(n)$ to denote 
the 2-qubit gate count for an $n$-qubit controlled-rotation gate in 
the form of $R = \sum^{2^{n-1}-1}_{i=0}R_{z}(\alpha_i) \otimes \outer{i}{i}$,  
we know that 
\begin{eqnarray}
\mathcal{U}(n) & = & 4 \times \mathcal{U}(n-1) + 3 \times \mathcal{R}(n) 
\label{eq:u-n-with-r} \\
\mathcal{U}(2) & = & 1    
\label{eq:u-2}
\end{eqnarray}

According to Lemma~\ref{rzcount}, we have that 
\begin{eqnarray}
\mathcal{U}(n) & = & 4 \times \mathcal{U}(n-1) + 3 \times (3 \times 2^{n-2} - 2)
\nonumber \\
& = & 4 \times \mathcal{U}(n-1) + 9 \times 2^{n-2} - 6
\label{eq:u-n-full-formula}
\end{eqnarray}

Now we can do an iterative calculation that uses
Equations~(\ref{eq:u-2})--(\ref{eq:u-n-full-formula}):
\begin{eqnarray}
\mathcal{U}(n) & = & 4 \times \mathcal{U}(n-1) + 9 \times 2^{n-2} - 6 \nonumber \\
& = & 4 \times (4 \times \mathcal{U}(n-2) + 9 \times 2^{n-3} - 6) + 9 \times 2^{n-2} - 6 \nonumber \\
& = & 4^2 \times \mathcal{U}(n-2) + 9 \times (4 \times 2^{n-3} + 1 \times 2^{n-2}) - 6 \times (4 + 1) \nonumber \\
&& \cdots \nonumber \\
& = &  4^{n-2} \times \mathcal{U}(2) + 9 \times \sum^{n-3}_{i=0}(4^i \times 2^{n-2-i}) - 6 \times \sum^{n-3}_{i=0} 4^i  \nonumber \\
& = &  4^{n-2}  + (9 \times 2^{n-2}) \times (2^{n-2} - 1) - 2 \times (4^{n-2} - 1) \nonumber \\
& = &  2 \times 4^{n-1} - 9 \times 2^{n-2} + 2 
\nonumber
\end{eqnarray}

We conclude that we can implement $U$ by 
$(\frac{1}{2} \times 4^n - 9 \times 2^{n-2} + 2 )$ 2-qubit gates. 


\paragraph{Transition to MP3:}
Now we can finish the proof by focusing on approximation.

\paragraph{Main Point 3:}
We will prove the full theorem by approximating the $R_z$ gates. \\
\mbox{} \\ \hspace*{1cm}
   [We define approximation of unitaries] \\
We define the Hilbert-Schmidt distance $d$ for 
$n$-qubit $n$-qubit unitaries $U, V$:
\(
d(U,V) &=& 
\sqrt{ 1 - \frac{ || Tr( U^\dagger V ) ||^2 }{ N^2 } }
\)
where $Tr(U)$ denotes the trace of $U$, and $N = 2^n$.

% We have that $d$ satisfies the triangle inequality:
% $d(U,W) \leq d(U,V) + d(V,W)$.

\newpage

\mbox{} \\ \hspace*{1cm}
   [We state some known lemmas about the Hilbert-Schmidt distance] % \\

\begin{lemma}
\label{lem:hs-distance-to-itself-is-zero}
$d(U,U) = 0$. \\
\end{lemma}

\begin{proof}
Notice $Tr( U^\dagger U ) = Tr( I_N ) = N$.
We calculate:
\(
d(U,U) &=& \sqrt{ 1 - \frac{ || Tr( U^\dagger U ) ||^2 }{ N^2 } } 
 \;\;=\;\; \sqrt{ 1 - \frac{ || N ||^2 }{ N^2 } } 
 \;\;=\;\; 0
\)
\end{proof}

\begin{lemma} {\bf (Wang and Zhang, 1994)}
\(
\sqrt{ 1 - \frac{ || Tr( UV ) ||^2 }{ N^2 } }
&\leq&
\sqrt{ 1 - \frac{ || Tr( U ) ||^2 }{ N^2 } }
\;+\;
\sqrt{ 1 - \frac{ || Tr( V ) ||^2 }{ N^2 } }
\)
\label{lem:trace-inequality}
\end{lemma}

For approximation of a product of gates, errors add at most linearly.

\begin{lemma}
$d(U_1 U_2, V_1 V_2) \;\leq\; d(U_1, V_1) \;+\; d(U_2, V_2)$.
\label{lem:hs-small-product-rule}
\end{lemma}

\begin{proof}
Notice that
$Tr((U_1 U_2)^\dagger\ (V_1 V_2)) \;=\;
Tr(U_2^\dagger U_1^\dagger V_1 V_2) 
\;=\;
Tr(U_1^\dagger V_1 V_2 U_2^\dagger)$.
From this we derive
\(
d(U_1 U_2, V_1 V_2) &=&
\sqrt{ 1 - \frac{ || Tr( (U_1 U_2)^\dagger\ (V_1 V_2) ) ||^2 }{ N^2 } }
\\ &=&
\sqrt{ 1 - \frac{ || Tr( (U_1^\dagger V_1)\ (V_2 U_2^\dagger) ) ||^2 }{ N^2 } }
\\ &\leq&
\sqrt{ 1 - \frac{ || Tr( U_1^\dagger V_1 ) ||^2 }{ N^2 } }
\;+\;
\sqrt{ 1 - \frac{ || Tr( V_2 U_2^\dagger ) ||^2 }{ N^2 } }
\\ &=&
\sqrt{ 1 - \frac{ || Tr( U_1^\dagger V_1 ) ||^2 }{ N^2 } }
\;+\;
\sqrt{ 1 - \frac{ || Tr( U_2^\dagger V_2 ) ||^2 }{ N^2 } }
\\[2mm] &=&
d(U_1,V_1) + d(U_2,V_2)
\)
In the first step, we use the definition of $d$.
In the second step, we use $(UV)^\dagger = V^\dagger\ U^\dagger$ and
$Tr(UV) = Tr(VU)$.
In the third step, we use Lemma~\ref{lem:trace-inequality}.
In the fourth step, we use $Tr(UV) = Tr(VU)$.
In the fifth step, we use the definition of $d$.
\end{proof}


\begin{lemma}
$d(\ U_1 \ldots U_n,\ \ V_1 \ldots V_n\ ) \;\leq\; \sum_{i=1}^n d(U_i,V_i)$.
\label{lem:hs-big-product-rule}
\end{lemma}

\begin{proof}
We proceed by induction on $n$.

In the base case of $n=1$, we have $d(U_1,V_1) \leq d(U_1,V_1)$, as required.

In the induction step, suppose the lemma holds for $n$.
Consider the case of $n+1$.
\(
d(\ U_1 \ldots U_{n+1},\ \ V_1 \ldots V_{n+1}\ ) 
&\leq&
d(\ U_1 \ldots U_n,\ \ V_1 \ldots V_n\ ) + d(U_{n+1},V_{n+1}) 
\\ &\leq&
\sum_{i=1}^n d(U_i,V_i) + d(U_{n+1},V_{n+1})
\\ &=&
\sum_{i=1}^{n+1} d(U_i,V_i) 
\)
In the first step, we use Lemma~\ref{lem:hs-small-product-rule}.
\end{proof}

\newpage

\mbox{} \\ \hspace*{1cm}
   [We prove the full version of the theorem] \\
We will rely on that any 2-qubit gate can be implemented using 
2 $C(X)$ gates and some 1-qubit gates in $\{H, T, R_z(\theta)\}$.

\begin{lemma} {\bf (Shende, Markov, and Bullock 2004)}
For any 2-qubit unitary $U$,
there exists a circuit $C$ built from 
$\{C(X), H, T, R_z(\theta)\}$, 
such that $U = C$.
\label{lem:implementation-of-a-2-qubit-gate}
\end{lemma}

We will rely on that $\{H, T\}$ can approximate $R_z(\theta)$:
we can use $\{H, T\}$ to get as close as we wish to any $R_z(\theta)$ gate.
Intuitively, $\{H, T\}$ is universal for defining $R_z(\theta)$ gates,
as expressed in Lemma~\ref{lem:approximation-of-rz}.

\begin{lemma}
For any $\epsilon > 0$,
there exists a circuit $C$ built from $\{H, T\}$,
such that $d(R_z(\theta),C) < \epsilon$.
\label{lem:approximation-of-rz}
\end{lemma}

Now we can prove the theorem, which we repeat here.

\bigskip

\noindent
{\bf Theorem}.  For any unitary $U$ and any $\epsilon > 0$, 
there exists a circuit $C$ built from $\{C(X), H, T\}$,
such that $d(U,C) < \epsilon$.

\begin{proof}
For a unitary $U$ and $\epsilon > 0$,
Lemma~\ref{lem:decomposition-to-an-easy-gate-set} says that
there exists a circuit $C_1$ built from 2-qubit gates and
$\{H, S, S^\dagger, R_z(\theta)\}$, 
such that $U = C_1$.
We have that in $C_1$
we can use Lemma~\ref{lem:implementation-of-a-2-qubit-gate} 
to replace every 2-qubit gate that is different from $C(X)$ with
a circuit built from $\{C(X), H, T, R_z(\theta)\}$.
We also have that in $C_1$ we can replace $S$ and $S^\dagger$ based on
the following equalities:
$S = T T$ and $S^\dagger = T T T T T T$.
This leaves us with a circuit $C_2$ built from 
$\{C(X), H, T, R_z(\theta)\}$ 
such that $U = C_1 = C_2$.

Let $k$ be the number of $R_z$ gates in $C_2$, and
let $\theta_1, \ldots, \theta_k$ be the angles used in those gates.
From Lemma~\ref{lem:approximation-of-rz} we get circuits $W_j$ such that
\begin{eqnarray}
d(R_z(\theta_j),W_j) &<& \epsilon/k
\label{eq:approximation-of-r-z}
\end{eqnarray}
Let $C_3$ be a version of $C_2$ in which we have replaced
each $R_z(\theta_j)$ by $W_j$.
We calculate:
\(
d(U, C_3) &=& 
d(C_2, C_3) \;\;\leq\;\;
\sum_{j=1}^k d(R_z(\theta_j),W_j) \;\;\leq\;\;
k \cdot (\epsilon/k) \;\;=\;\; 
\epsilon
\)
In the second step, we use 
Lemma~\ref{lem:hs-distance-to-itself-is-zero} and
Lemma~\ref{lem:hs-big-product-rule}.
In the third step, we use Equation~(\ref{eq:approximation-of-r-z}).
\end{proof}


\paragraph{Transition to Close:} 
So there you have it.

\paragraph{Review:}
The key property of a universal gate set is that 
it can approximate any unitary.
Our proof strategy is to first prove the theorem for an easy gate set and then
extend to $\{C(X), H, T\}$.

\paragraph{Strong finish:}
Quantum computing has probabilities built in and
the idea of a universal gate set has approximation built in.
Quantum algorithms have no requirement to be exact; 
rather, they can try to get close.

\paragraph{Call to action:}
Implement the procedure for compiling any unitary into the easy gate set.

\end{document}

